\documentclass{situationist-preprint}

\title{Inductive Risk}
\msid{SCI-13-b}
\msfamily{SCI-13}
\msdiscipline{Philosophy of Science}
\mstype{Consequence}
\mskeywords{inductive risk; values in science; acceptance; thresholds; Rudner; Douglas; wishful thinking}

\begin{document}
\maketitle

\begin{abstract}
SCI-13-a holds that non-epistemic values legitimately set the threshold for accepting a hypothesis but do not determine the content of accepted results. We derive what this implies for the management of inductive risk. Values may set different thresholds for different hypotheses within one inquiry and may enter through methodological choices that function as thresholds; they may weigh the costs of error but not the desirability of a hypothesis's being true; and, because the evidential report is common to all inquirers, any user of a result can redetermine its acceptance at her own threshold. The argument from inductive risk thus licenses a definite and limited role for values, and requires a practice of reporting that makes that role transparent.
\end{abstract}

\section{The argument and its scope}

\citet{Rudner1953} argued that since no hypothesis is verified with certainty, accepting one requires deciding whether the evidence is strong enough, and that this decision depends on how serious it would be to be wrong. \citet{Jeffrey1956} replied that scientists should report probabilities rather than accept hypotheses; \citet{Douglas2000,Douglas2009} answered that the choice of methods and the characterization of evidence already involve decisions under risk. SCI-13-a accepts Rudner's premise in the form of its Principle~P1: the threshold for acceptance may depend on the costs of error, including non-epistemic costs. Its Principle~P2 requires the content of an accepted result, the hypothesis and its evidential report, to be fixed by the evidence. Theorem~3.3 of SCI-13-a concludes that values can decide whether a result is accepted, not what it says. We draw four consequences.

\section{Consequences}

\begin{corollary}[Differential thresholds]\label{c:diff}
Within a single inquiry, different hypotheses may legitimately be accepted at different thresholds, and a hypothesis and its negation may have different thresholds.
\end{corollary}

\begin{proof}
By Principle~P1 of SCI-13-a, the threshold depends on the costs of wrongly accepting the hypothesis. These costs differ between hypotheses and between a hypothesis and its negation: wrongly accepting that a compound is safe differs in its consequences from wrongly accepting that it is toxic.
\end{proof}

Corollary~\ref{c:diff} is compatible with Proposition~3.1 of SCI-13-a. Since every threshold exceeds one half (Definition~2.2 of SCI-13-a), a hypothesis and its negation are never both accepted on the same evidence, whatever their thresholds.

\begin{corollary}[Methodological choices]\label{c:method}
A methodological choice whose effect is to fix how much evidence is required before a result is accepted, such as a significance level or a stopping rule, may legitimately be guided by non-epistemic values, provided that it is reported.
\end{corollary}

\begin{proof}
Such a choice is a threshold setting in the sense of Principle~P1 of SCI-13-a. The proviso follows from Principle~P2, which requires the report to state the methods that produced the evidence.
\end{proof}

\begin{corollary}[No wishful thinking]\label{c:wish}
Non-epistemic values may bear on the threshold only through the costs of error. A preference for a hypothesis's being true is not a legitimate ground for lowering its threshold.
\end{corollary}

\begin{proof}
A threshold lowered because the hypothesis is desirable, rather than because wrongly accepting it is cheap, makes which hypothesis is accepted depend on which hypothesis is wanted. That is a determination of content by value, which Principle~P2 of SCI-13-a excludes.
\end{proof}

Corollary~\ref{c:wish} is the line Douglas draws between the indirect role of values, in weighing the consequences of error, and their illegitimate direct role, as reasons for belief \citep{Douglas2009}. It follows here from the content principle rather than being added to it.

\begin{corollary}[Recomputability]\label{c:recompute}
Given a published evidential report, any user of a result can determine whether the hypothesis would be accepted at her own threshold.
\end{corollary}

\begin{proof}
By Proposition~3.2 of SCI-13-a, inquirers who share their evidence issue the same report, and by Principle~P2 the report states the degree of support. Acceptance at threshold $\tau$ is a comparison of that support with $\tau$ (Definition~2.2 of SCI-13-a).
\end{proof}

Corollary~\ref{c:recompute} answers the worry that the argument from inductive risk lets the values of scientists be imposed on the public. Whatever threshold a scientist used, a regulator or a citizen with different values can reach her own verdict from the same report. This vindicates Jeffrey's emphasis on reporting without his conclusion that scientists should not accept hypotheses. Acceptance remains necessary, because results must enter the record as premises for further work; the report keeps it revisable by those whose costs of error differ.

\begin{objection}
Values enter not only at the threshold but in characterizing the data, for instance in classifying ambiguous tissue samples as malignant or benign \citep{Douglas2000}. Here values seem to shape content.
\end{objection}

\begin{reply}
The classification of an ambiguous sample is the acceptance of a hypothesis about that sample, and the inquirer who classifies conservatively is setting a threshold. By Corollary~\ref{c:method}, the practice is legitimate if it is reported; by Corollary~\ref{c:recompute}, a reader with other values can reclassify. What the values do not do is change what the evidence about each sample is.
\end{reply}

\section{Conclusion}

Inductive risk is managed by thresholds, and thresholds are the only point at which values enter. They may vary across hypotheses and enter through method, but they weigh the costs of error, not the attractions of a conclusion, and the common report lets anyone reset them.

\begin{thebibliography}{99}\small
\bibitem[Douglas(2000)]{Douglas2000} Douglas, H. (2000). Inductive risk and values in science. \emph{Philosophy of Science}, 67(4), 559--579.
\bibitem[Douglas(2009)]{Douglas2009} Douglas, H. (2009). \emph{Science, Policy, and the Value-Free Ideal}. University of Pittsburgh Press.
\bibitem[Jeffrey(1956)]{Jeffrey1956} Jeffrey, R.~C. (1956). Valuation and acceptance of scientific hypotheses. \emph{Philosophy of Science}, 23(3), 237--246.
\bibitem[Rudner(1953)]{Rudner1953} Rudner, R. (1953). The scientist \emph{qua} scientist makes value judgments. \emph{Philosophy of Science}, 20(1), 1--6.
\end{thebibliography}

\end{document}
